\documentclass[12pt,a4paper]{article}

\usepackage[utf8]{inputenc}
\usepackage[T1]{fontenc}
\usepackage[english]{babel}
\usepackage{geometry}
\usepackage{booktabs}
\usepackage{longtable}
\usepackage{array}
\usepackage{tabularx}
\usepackage{enumitem}
\usepackage{graphicx}
\usepackage{xcolor}
\usepackage{hyperref}
\usepackage{fancyhdr}

\geometry{a4paper, left=2.5cm, right=2.5cm, top=2.5cm, bottom=2.5cm}
\setlength{\parindent}{1.25cm}
\setlength{\parskip}{0.25em}
\setlist{itemsep=0.2em, topsep=0.3em}
\hypersetup{colorlinks=true, linkcolor=blue, urlcolor=blue}

\pagestyle{fancy}
\fancyhf{}
\lhead{Stock Market Project}
\rhead{Documentation}
\cfoot{\thepage}

\newcommand{\code}[1]{\texttt{\small #1}}
\newcommand{\class}[1]{\textbf{\code{#1}}}

\begin{document}

\begin{titlepage}
\centering
\vspace*{\fill}

{\LARGE\bfseries Documentation\\[0.3cm]
for the Stock Market Project\par}

\vspace{0.7cm}
{\large Distributed system for stock-market simulation and portfolio management\par}

\vspace{1.5cm}
{\large Student: Vladu Alexandru-Gheorghe\par}
\vspace{0.4cm}
{\large Group: 244\par}

\vspace{1cm}
{\large Academic year: 2026-2027\par}

\vspace*{\fill}
\end{titlepage}

\tableofcontents
\newpage

\section{Purpose and general description}

The project implements a minimal distributed system for a stock market. It
consists of two independent Java web services connected through Apache Kafka:

\begin{itemize}
    \item \emph{stock-market-service} simulates the market, maintains the
    initial list of companies, and publishes prices to the Kafka topic
    \code{stock-prices};
    \item \emph{asset-management-service} consumes price updates, stores the
    latest quote for each symbol, and provides a web interface for viewing and
    managing a portfolio.
\end{itemize}

The producer and consumer are decoupled through Kafka, which stores and
delivers events independently of the web interface.

The solution uses Java 21, Jakarta Enterprise Beans (EJB), Jakarta Servlet,
JSP, and JDBC connection pools. Jetty 12 hosts the producer service and
TomEE Plus 10.2 hosts the asset-management service. Both applications are packaged
as WAR archives. Business data is stored in PostgreSQL, Microsoft SQL Server,
and Oracle Database; Kafka messages use a compact project-defined JSON
representation.

\section{System architecture}

\subsection{Components}

The logical architecture is:

\begin{figure}[htbp]
    \centering
    \includegraphics[width=\textwidth]{architecture.png}
    \caption{Simplified system architecture, including services, Kafka,
    database connections, and database migrations}
    \label{fig:architecture}
\end{figure}

\noindent
The browser uses the TomEE-hosted asset-management service. The Jetty service
publishes prices to Kafka; the consumer reads them and uses one pool per
database. The migration service creates the PostgreSQL, SQL Server, and
OracleDB schemas.

\begin{center}
\begin{tabular}{>{\raggedright\arraybackslash}p{3.3cm}
                >{\centering\arraybackslash}p{1.5cm}
                >{\raggedright\arraybackslash}p{7.0cm}}
\toprule
\textbf{Component} & \textbf{Port} & \textbf{Responsibility}\\
\midrule
Kafka & 9092 & Transports quote events through the \code{stock-prices} topic.\\
Stock market service & 8082 & Runs on Jetty, generates prices, and publishes
messages.\\
Asset-management service & 8081 & Runs on TomEE, consumes messages, and
exposes the web API and JSP interface.\\
PostgreSQL & 5432 & Stores users and authentication data.\\
Microsoft SQL Server & 1433 & Stores portfolios and cash balances.\\
OracleDB & 1521 & Stores the trade ledger.\\
\bottomrule
\end{tabular}
\end{center}

At startup, the market listener creates the repository, publisher, and
simulator. It publishes AAPL, AMZN, MSFT, NVDA, and TSLA, then schedules
updates every 5--15 seconds.

At startup, EJB components create the configuration and three HikariCP pools.
The consumer reads the price topic and stores the latest quote and timestamp.
Repositories use PostgreSQL for users, SQL Server for portfolios, and Oracle
for trades.

\subsection{Database model and connections}

The databases are deliberately separated by business responsibility. Their
tables are different, but the application joins them through stable
identifiers because a relational foreign key cannot directly span three
database engines:

\begin{itemize}
    \item PostgreSQL contains \code{users(user\_id, username, email,
    created\_at)} and authentication records.
    \item Microsoft SQL Server contains
    \code{portfolios(portfolio\_id, user\_id, portfolio\_name,
    cash\_balance, created\_at)}. \code{user\_id} is a logical reference to
    PostgreSQL.
    \item OracleDB contains
    \code{trades(trade\_id, portfolio\_id, symbol, side, quantity, price,
    executed\_at)}. \code{portfolio\_id} is a logical reference to SQL Server.
\end{itemize}

The connections use \code{POSTGRES\_URL}, \code{SQLSERVER\_URL}, and
\code{ORACLE\_URL}, with separate credentials. The singleton
\class{Database} owns the HikariCP pools; repositories use prepared
statements and try-with-resources. Flyway applies each database's migrations.

The conceptual JPA model has one persistence unit per database:
\class{UserEntity}, \class{PortfolioEntity}, and \class{TradeEntity}. Since
JPA cannot create cross-database foreign keys, associations use identifier
fields. The current repositories use JDBC to keep vendor-specific connections
explicit.

\subsection{Processing sequence}

\begin{enumerate}
    \item Docker Compose starts Kafka and the databases; health checks confirm
    availability.
    \item Maven builds and deploys both WAR archives.
    \item The market-service listener publishes the initial stock list.
    \item The consumer converts each event to a \class{StockCompany}, stores
    it by symbol, and notifies SSE listeners.
    \item The web client requests pages, JSON, or an SSE stream.
    \item Trading conditionally updates SQL Server cash and inserts the Oracle
    trade; the dashboard joins records from all three databases.
\end{enumerate}

\subsection{Concurrency and consistency}

Concurrent collections protect quotes and listeners. \class{TradingControl}
uses local JDBC transactions and compensating actions: a failed Oracle insert
reverses the SQL Server cash movement. EJB concurrency controls and database
constraints prevent conflicting trades and negative balances.

\section{Inputs, outputs, and operating conditions}

\subsection{Configuration and dependencies}

The system requires JDK 21, Maven, Docker with Docker Compose, PostgreSQL 16,
Microsoft SQL Server 2022, Oracle Database Free 23, Jetty 12, and TomEE Plus 10.2.
Docker Compose installs and starts the three databases on the local computer.
The scripts assume a Homebrew-based environment for locating the application
servers. For a different installation, the variables
\code{CATALINA\_HOME}, \code{CATALINA\_BASE}, \code{JETTY\_HOME}, and
\code{JETTY\_BASE} can be changed.

The supported environment variables are:

\begin{longtable}{@{}>{\raggedright\arraybackslash}p{0.27\linewidth}
                  >{\raggedright\arraybackslash}p{0.36\linewidth}
                  >{\raggedright\arraybackslash}p{0.29\linewidth}@{}}
\toprule
\textbf{Variable} & \textbf{Default} & \textbf{Purpose}\\
\midrule
\endhead
\code{KAFKA\_\allowbreak BOOTSTRAP\_\allowbreak SERVERS} & \code{localhost:9092} & Kafka broker address.\\
\code{KAFKA\_TOPIC} & \code{stock-prices} & Event topic.\\
\code{POSTGRES\_URL} &
\code{jdbc:\allowbreak postgresql://\allowbreak localhost:\allowbreak
5432/\allowbreak stock\_market}
  & PostgreSQL connection.\\
\code{SQLSERVER\_URL} &
\code{jdbc:\allowbreak sqlserver://\allowbreak localhost:\allowbreak
1433;\allowbreak databaseName=stock\_market;\allowbreak encrypt=false}
  & Microsoft SQL Server connection.\\
\code{ORACLE\_URL} &
\code{jdbc:\allowbreak oracle:\allowbreak thin:@\allowbreak localhost:\allowbreak
1521/\allowbreak FREEPDB1}
  & OracleDB connection.\\
\code{POSTGRES\_PASSWORD}, \code{SQLSERVER\_PASSWORD},
\code{ORACLE\_APP\_PASSWORD} & required & Passwords for the three pools.\\
\bottomrule
\end{longtable}

A market message contains \code{symbol}, \code{name}, \code{price}, and
\code{quantity}. The symbol and name must not be blank, the price must be a
non-negative decimal value, and the quantity must be a non-negative integer.
Transaction quantities may be fractional, but must be strictly positive.

\subsection{Build and deployment}

Each service is packaged as a WAR with Maven and deployed under the root
context by its deployment script. The asset-management WAR runs on TomEE
port 8081; the stock-market WAR runs on Jetty port 8082.

The complete local deployment can be started from the project root with:

\begin{verbatim}
./run.sh
\end{verbatim}

The script starts Kafka and the databases, waits for health checks, applies
migrations, deploys both WARs, and starts the servers. Logs are written to
\code{/tmp/asset-management-service-tomcat.log} and
\code{/tmp/stock-market-service-jetty.log}.

Deployment requires Docker to be running and assumes that the server
locations are available through Homebrew. The database credentials must be
provided through \code{.env} or the shell environment. The paths can be adapted by changing
\code{CATALINA\_HOME}, \code{CATALINA\_BASE}, \code{JETTY\_HOME}, and
\code{JETTY\_BASE}. To stop the local deployment from another terminal, use:

\begin{verbatim}
./stop.sh
\end{verbatim}

This releases ports 8081 and 8082; \code{Ctrl+C} also stops the servers and
Docker Compose stack.

\subsection{HTTP interface}

\begin{longtable}{>{\raggedright\arraybackslash}p{2.2cm}
                  >{\raggedright\arraybackslash}p{4.1cm}
                  >{\raggedright\arraybackslash}p{7.2cm}}
\toprule
\textbf{Method} & \textbf{Route} & \textbf{Behavior}\\
\midrule
\endhead
GET & \code{/} & Application home page.\\
GET & \code{/assets} & JSP page containing portfolio data.\\
GET & \code{/assets/events} & SSE stream with updated market and account data.\\
GET & \code{/live-market} & Live-market user interface.\\
GET & \code{/live-market/data} & JSON containing stocks and account status.\\
GET & \code{/live-market/events} & SSE stream for live updates.\\
POST & \code{/live-market/buy} & Buys the requested quantity of a symbol.\\
POST & \code{/live-market/sell} & Sells the requested quantity of a symbol.\\
GET/POST & \code{/api/users} & User CRUD operations in PostgreSQL.\\
GET/POST & \code{/api/portfolios} & Portfolio CRUD operations in SQL Server.\\
GET/POST & \code{/api/trades} & Trade CRUD operations in OracleDB.\\
GET & \code{/api/dashboard} & Joined data from all three databases.\\
\bottomrule
\end{longtable}

CRUD routes receive resource-specific JSON. Market operations receive
\code{symbol} and \code{quantity}; invalid input, unknown stocks, insufficient
cash, or insufficient holdings produce HTTP 400 with an \code{error} field.
The dashboard joins user, portfolio, and trade data in the application.

\section{Objects and prototypes}

The following sections describe the project objects without including their
implementation. The prototypes show method names, parameters, and return
types.

\subsection{Stock-market service}

\begin{description}
\item[\class{StockCompany}] A record with the components \code{String symbol},
\code{String name}, \code{BigDecimal price}, and \code{long quantity}. The
record's implicit constructor creates a complete quote.

\item[\class{StockCompanyJson}] A utility class with no public instantiation.
\begin{itemize}
    \item \code{static String write(StockCompany company)} -- serializes one
    quote.
    \item \code{static String writeAll(List<StockCompany> companies)} --
    serializes a list.
    \item \code{static StockCompany fromParameters(String symbol, String name,
    String price, String quantity)} -- validates text parameters and creates a
    quote.
\end{itemize}
The parsing method can throw \code{IllegalArgumentException} for missing
fields, non-numeric values, or negative values.

\item[\class{StockCompanyRepository}] The no-argument constructor loads five
initial companies.
\begin{itemize}
    \item \code{void save(StockCompany company)} -- stores a company by
    \code{symbol}.
    \item \code{List<StockCompany> findAll()} -- returns companies in sorted
    order.
    \item \code{StockCompany createRandomPriceUpdate()} -- selects a company
    and calculates a random price variation.
\end{itemize}

\item[\class{StockMarketPublisher}] The no-argument constructor configures the
Kafka producer. \code{RecordMetadata publish(StockCompany company)} publishes
the message and waits for broker confirmation; it can throw \code{Exception}.
\code{void close()} closes the producer and implements \code{AutoCloseable}.

\item[\class{StockMarketSimulator}] The constructor
\code{StockMarketSimulator(StockCompanyRepository repository,
StockMarketPublisher publisher)} receives the data source and publisher.
\code{void start()} publishes initial values and starts scheduling.
\code{void close()} cancels the next update and stops the executor.
\end{description}

\subsection{Asset-management service}

\begin{description}
\item[\class{StockCompany}] A DTO record with \code{String symbol},
\code{String name}, \code{BigDecimal price}, and \code{long quantity}. It is
the local representation of a consumed message.

\item[\class{StockCompanyJson}] A utility class without a public constructor.
\code{static StockCompany read(String value)} parses a JSON message and can
throw \code{IllegalArgumentException} for an invalid format.

\item[\class{StockPriceConsumer}] The no-argument constructor configures the
Kafka consumer. Its public methods are:
\begin{itemize}
    \item \code{void start()} -- starts the consumer thread;
    \item \code{List<StockCompany> latestPrices()} -- returns sorted quotes;
    \item \code{Instant lastUpdated(String symbol)} -- returns an update time;
    \item \code{Map<String, StockCompany> latestPricesBySymbol()} -- returns
    the market;
    \item \code{void addListener(...)} and \code{void removeListener(...)} --
    manage SSE subscribers;
    \item \code{void close()} -- stops the consumer.
\end{itemize}

\item[\class{Database}] A singleton EJB that creates three HikariCP pools.
\code{DataSource postgres()}, \code{DataSource sqlServer()}, and
\code{DataSource oracle()} return the corresponding data sources.

\item[\class{JdbcUserRepository}] A stateless EJB for PostgreSQL. Its methods
\code{User find(long userId)}, \code{List<User> findAll()},
\code{User create(User user)}, \code{User update(User user)}, and
\code{void delete(long userId)} implement user CRUD and can throw
\code{SQLException} or \code{NotFoundException}.

\item[\class{JdbcPortfolioRepository}] A stateless EJB for SQL Server. Its
methods \code{Portfolio find(long portfolioId)}, \code{List<Portfolio>
findByUser(long userId)}, \code{Portfolio create(Portfolio portfolio)},
\code{Portfolio update(Portfolio portfolio)}, and \code{void delete(long
portfolioId)} implement portfolio CRUD.

\item[\class{JdbcTradeRepository}] A stateless EJB for OracleDB. Its methods
\code{Trade find(long tradeId)}, \code{List<Trade> findByPortfolio(long
portfolioId)}, \code{Trade create(Trade trade)}, \code{Trade update(Trade
trade)}, and \code{void delete(long tradeId)} implement trade CRUD. Trade
updates and deletes are audited and validated against portfolio ownership.

\item[\class{UserControl}], \class{PortfolioControl}, and
\class{TradingControl} are EJB services for validation, ownership,
repository calls, and cross-database coordination.
\code{DashboardControl} joins the three repository results into a dashboard
DTO and reports failures explicitly.

\item[\class{UserServlet}], \class{PortfolioServlet},
\class{TradeServlet}, and \class{DashboardServlet} extend
\code{HttpServlet} and expose CRUD or joined dashboard data as JSON.
\Their HTTP methods can throw \code{ServletException} and
\code{IOException}; validation errors become HTTP 400 responses.

\item[\class{AssetServlet}] and \class{LiveMarketServlet} extend
\code{HttpServlet} and serve JSP pages, JSON data, and SSE updates.
\Their \code{doGet} and \code{doPost} methods delegate business logic to EJB
controls.
\end{description}

\class{StockMarketContextListener} and \class{StockPriceContextListener}
implement \code{ServletContextListener}; their initialization and destruction
methods create, register, and release service resources.

\section{System calls and relevant libraries}

System-level operations are provided by Java libraries and launch scripts:

\begin{itemize}
    \item TCP/IP through the Kafka client and the Jetty/TomEE HTTP servers;
    \item Java threads for the consumer and scheduled simulator;
    \item JDBC \code{DataSource}, \code{Connection}, and prepared statements;
    \item JPA entities and persistence units as the conceptual ORM model;
    \item Docker Compose, Maven, TomEE, and Jetty started by scripts;
    \item Kafka wake-up and executor interruption during shutdown;
    \item persistent SSE connections using \code{text/event-stream}.
\end{itemize}

\section{Usage examples}

\subsection{Starting the system}

From the project directory, run:

\begin{verbatim}
./run.sh
\end{verbatim}

The script starts the infrastructure, deploys both services, and starts the
servers. The application is available at
\url{http://localhost:8081/live-market}; stop it with \code{Ctrl+C} or
\code{./stop.sh}.

\subsection{Buying and selling}

A client can buy 0.5 AAPL shares with:

\begin{verbatim}
curl -X POST http://localhost:8081/live-market/buy \
  -d "symbol=AAPL&quantity=0.5"
\end{verbatim}

The response is JSON containing the available balance, invested value, total
value, and holdings. A partial sale is analogous:

\begin{verbatim}
curl -X POST http://localhost:8081/live-market/sell \
  -d "symbol=AAPL&quantity=0.2"
\end{verbatim}

The current state can be read without modifying the databases:

\begin{verbatim}
curl http://localhost:8081/live-market/data
\end{verbatim}

The \code{/live-market} page uses SSE for live prices and values the account
using the latest server-side price.

\subsection{CRUD and cross-database examples}

A user is created in PostgreSQL, then a portfolio referencing its
\code{user\_id} is created in SQL Server:

\begin{verbatim}
curl -X POST http://localhost:8081/api/users \
  -H 'Content-Type: application/json' \
  -d '{"username":"alex","email":"alex@example.com"}'
curl -X POST http://localhost:8081/api/portfolios \
  -H 'Content-Type: application/json' \
  -d '{"userId":1,"portfolioName":"Main","cashBalance":10000}'
\end{verbatim}

A trade is then inserted in Oracle using the returned \code{portfolioId}.
The dashboard combines all three records:

\begin{verbatim}
curl http://localhost:8081/api/dashboard?userId=1
\end{verbatim}

The same resources support \code{GET}, \code{PUT}, and \code{DELETE}. Before
creating or updating a child record, the EJB service verifies the parent in
the preceding database. Deletion is performed in reverse dependency order:
trades, portfolios, then users.

\section{Alternative solutions and limitations}

Kafka could be replaced by REST, WebSocket, or JMS. REST is simpler for a
single request but couples the services; WebSocket is suitable for browser
updates but does not replace broker persistence. SSE is used because browser
updates are unidirectional.

The implementation has one local broker, one instance of each database, and
compensating rather than distributed transactions. Production would require
secret management, replication, stronger validation, structured logging, and
an XA-capable strategy. JPA could replace the JDBC repositories, but the
three persistence units would still need explicit coordination.

\section{Conclusions}

The project demonstrates a clear distributed architecture in which market-data
generation, message transport, and portfolio management are separate
responsibilities. Kafka provides event decoupling and delivery, EJBs
coordinate application services and repositories, JDBC/JPA define the
persistence boundary, and Servlets/JSP provide the web interface. PostgreSQL,
SQL Server, and Oracle store different but related tables, while the
dashboard performs the cross-database join in the application. Through the
documented CRUD routes, the behavior can be verified both component by
component and end to end.

\end{document}
